\chapter{Introducción}

Este documento contiene el \textbf{estudio de conexión simplificado} del proyecto solar fotovoltaico que se instalará en el supermercado MAS X MENOS sede Guarín (Bucaramanga, Santander). El análisis evalúa el impacto de la generación sobre el Sistema de Distribución Local (SDL) de ESSA, específicamente sobre el \textbf{circuito 10-502} alimentado desde la subestación Conucos a 13,8~kV.

El estudio se elaboró a partir de la información oficial de red entregada por el operador (archivo \textit{Datos Circuito 10 502 - 2026.xlsx} y diagrama unifilar \textit{CTO 10 502.pdf}), e incluye análisis de flujo de carga, cortocircuito, protecciones y pérdidas técnicas, conforme a la Resolución CREG 030 de 2018 y demás normativa aplicable.

\section{Resumen Ejecutivo}

El supermercado MAS X MENOS sede Guarín contempla la instalación de un sistema fotovoltaico de \textbf{141,6~kWp} (DC) / \textbf{120~kW} (AC), compuesto por 240 paneles JINKO de 590~W y dos inversores SOLIS 60K-LV-5G de 60~kW. La relación DC/AC es de 1,18.

La conexión se realizará en el \textbf{nodo 3272966} del circuito 10-502, mediante un transformador de 150~kVA, 13,2~kV/220~V, grupo Dyn5. La carga del cliente modelada en dicho nodo es de 150~kVA con factor de potencia 0,9 (135~kW).

\begin{table}[H]
    \centering
    \caption{Datos clave del proyecto y del punto de conexión}
    \begin{tabular}{ll}
    \toprule
    \textbf{Parámetro} & \textbf{Valor} \\
    \midrule
    Proyecto & Solar MAS X MENOS Guarín \\
    Capacidad DC / AC & 141,6 kWp / 120 kW \\
    Operador de red & ESSA \\
    Subestación / circuito & Conucos / 10-502 \\
    Nivel de tensión & 13,8 kV \\
    Punto de conexión (POC) & Nodo 3272966 \\
    Tramo de conexión & Línea 807675 (3272869--3272966A) \\
    Conductor de conexión & 3F\_15\_CU\_2\_XLPE, 3 m, $I_n=155$~A \\
    Demanda máxima del circuito & 1,683 MVA (12:00 h, 19/03/2024) \\
    Equivalente SC trifásico (SE) & 654,49 MVA / 27,38 kA \\
    \bottomrule
    \end{tabular}
\end{table}

De los resultados se destaca que, para la alternativa de conexión evaluada:

\begin{itemize}
    \item No se presentan violaciones a los límites de tensión (0,90--1,10~p.u.).
    \item La cargabilidad de líneas se reduce (consumo local de la generación); el transformador de conexión permanece por debajo del 56\%.
    \item El aporte de cortocircuito de los inversores es despreciable frente al equivalente de red; no se requieren cambios en las protecciones del circuito.
    \item Se observa una reducción de pérdidas técnicas de hasta 0,34~kW en las horas de análisis.
\end{itemize}
